\documentclass[11pt,a4paper]{article}
\usepackage[margin=1.1in]{geometry}
\usepackage{amsmath,amssymb,amsthm}
\usepackage{lmodern}
\usepackage[T1]{fontenc}
\usepackage[hidelinks,pdfencoding=auto]{hyperref}
\hypersetup{
  pdftitle={An Erd\H os--Kac law for palindromes and for reversed primes},
  pdfauthor={Anonymous},
  pdfsubject={Analytic number theory},
  pdfkeywords={Erd\H os--Kac theorem; palindromes; reversed primes; sieve methods;
    level of distribution}}
\ifdefined\pdfgentounicode
  \input{glyphtounicode}
  \pdfgentounicode=1
\fi

\theoremstyle{plain}
\newtheorem{theorem}{Theorem}[section]
\newtheorem{proposition}[theorem]{Proposition}
\newtheorem{lemma}[theorem]{Lemma}
\newtheorem{corollary}[theorem]{Corollary}
\theoremstyle{definition}
\newtheorem{remark}[theorem]{Remark}
\newtheorem{example}[theorem]{Example}

\DeclareMathOperator{\Gal}{Gal}
\newcommand{\eps}{\varepsilon}
\numberwithin{equation}{section}

\title{An Erd\H os--Kac law for palindromes and for reversed primes}
\author{}
\date{}

\begin{document}
\maketitle

\begin{abstract}\noindent
For every integer $b\ge2$ we determine the distribution of the number of prime factors, counted with
or without multiplicity, of the base-$b$ palindromes with $\lambda$ digits and of the base-$b$
reversals of the primes with $\lambda$ digits: both obey an Erd\H os--Kac law with the classical
centring $\log\log b^{\lambda}$ and scaling $\sqrt{\log\log b^{\lambda}}$. For the number of distinct
prime factors we obtain in addition all moments of order up to $\tfrac12(\log\log b^{\lambda})^{1/3}$,
uniformly in the order. For reversed primes the law persists when the leading digit of the
prime is prescribed. The arithmetic inputs are a theorem of Col for palindromes and the
Bombieri--Vinogradov theorem of Dartyge, Rivat and Swaenepoel for reversed primes; in both cases the
excluded primes are those dividing $b(b^{2}-1)$. The deduction rests on the sieve-theoretic moment
estimate of Granville and Soundararajan, cast as a general criterion in the manner of Murty, Murty and
Pujahari and modified to allow a finite set of primes at which the local densities are arbitrary.
\end{abstract}

\medskip
\noindent{\small\emph{2020 Mathematics Subject Classification.} Primary 11N60; Secondary 11N36, 11A63.\\
\emph{Keywords.} Erd\H os--Kac theorem; palindromes; reversed primes; sieve methods; level of
distribution.}

\section{Introduction}

Let $\omega(n)$ be the number of distinct prime factors of $n$ and $\Omega(n)$ the number of prime
factors of $n$ counted with multiplicity. The theorem of Erd\H os and Kac \cite{EK} states that,
uniformly for $t\in\mathbb R$,
\begin{equation}\label{eq:EK}
\frac1x\,\#\Bigl\{n\le x:\ \omega(n)\le \log\log x+t\sqrt{\log\log x}\Bigr\}\ \longrightarrow\ \Phi(t)
\qquad(x\to\infty),
\end{equation}
where $\Phi$ is the standard normal distribution function. For sequences defined by conditions on the
digits, laws of this kind are known in several cases. Mauduit and S\'ark\"ozy proved an
Erd\H os--Kac theorem for the integers whose digit sum lies in a fixed residue class \cite{MS1}, and
for those whose digit sum is fixed \cite[Thm.~4]{MS2}. Erd\H os, Mauduit and S\'ark\"ozy
\cite[Thm.~1]{EMS} obtained an Erd\H os--Kac law for $\omega$ on the integers whose base-$b$ digits
are confined to a set $\mathcal D=\{0,d_2,\dots,d_s\}$ with $\gcd(d_2,\dots,d_s)=1$ and $s>\sqrt b$.
Konyagin, Mauduit and S\'ark\"ozy \cite{KMS} studied the normal order and the large values of $\omega$
on missing-digit sets, together with the large values of $\Omega$ on the integers with a fixed digit
sum. For palindromes and for reversals of primes we are not aware of such a law.

Fix an integer $b\ge2$. Every integer $n\ge0$ has a unique base-$b$ expansion
$n=\sum_{j\ge0}\eps_j(n)b^{j}$ with digits $\eps_j(n)\in\{0,1,\dots,b-1\}$, and we write
\begin{equation}\label{eq:defR}
R_\lambda(n)\ :=\ \sum_{j=0}^{\lambda-1}\eps_j(n)\,b^{\lambda-1-j}
\end{equation}
for the reversal of the $\lambda$ least significant digits of $n$. Let
\[
\mathcal T_\lambda\ :=\ \bigl\{n:\ b^{\lambda-1}\le n<b^{\lambda},\ \ n=R_\lambda(n)\bigr\}
\]
be the set of base-$b$ palindromes with exactly $\lambda$ digits, so that
$\#\mathcal T_\lambda=(b-1)b^{\lceil\lambda/2\rceil-1}$; and let
\[
\mathcal P_\lambda:=\bigl\{p\ \text{prime}:\ b^{\lambda-1}\le p<b^{\lambda}\bigr\},\qquad
\pi_\lambda:=\#\mathcal P_\lambda,
\]
so that the integers $R_\lambda(p)$ with $p\in\mathcal P_\lambda$ are the \emph{reversed primes}. The
term refers to the reversals themselves, which are in general composite; it is not to be confused with
the \emph{reversible primes} of \cite{DMRSS}, the primes whose reversal is again prime. This note
determines the distribution of $\omega$ and of $\Omega$ on both families, the palindromes and the
reversed primes.

For palindromes the question is old. Banks and Shparlinski \cite{BSh}, in their work on prime divisors
of palindromes, showed that for every large $\lambda$ some $n\in\mathcal T_\lambda$ has at least
$(\log\log n)^{1+o(1)}$ distinct prime factors. Their result concerns the largest value of $\omega$ on
$\mathcal T_\lambda$, and it does not address the typical value. We show in Theorem~\ref{thm:pal} below
that almost all palindromes, and not merely one, have $(1+o(1))\log\log n$ distinct prime factors, the
fluctuation around $\log\log n$ being of order $\sqrt{\log\log n}$. Concerning the typical order of
$\omega$ on $\mathcal T_\lambda$ nothing seems to be known, although the factorisation of palindromes
has been studied from other directions: Col \cite{Col} and Tuxanidy and Panario \cite{TP} produce palindromes
with boundedly many prime factors, Johnston and Kerr \cite{JK} show that infinitely many palindromes are
squarefree in every base, with an asymptotic count, and Chourasiya and Johnston \cite{CJ} obtain
cubefree palindromes in every base, and squarefree reversals of primes for $b\ge26000$.

Reversed primes have become accessible only recently. Bhowmik and
Suzuki \cite{BS1,BS2} established an effective Siegel--Walfisz theorem for the reversals $R_\lambda(p)$ in
arithmetic progressions, first for $b\ge31699$ and then for every $b\ge2$; Harm and Johnston \cite{HJ}
have since given a version in which the number of digits is not fixed. Dartyge, Rivat and
Swaenepoel \cite{DRS} obtained a Bombieri--Vinogradov theorem for reversals whose moduli range up to a
fixed power of $b^{\lambda}$; it is quoted below as Theorem~\ref{thm:drs}. Their own application is to a
sieve problem: they produce an explicit $\kappa_b$ such that $\Omega(R_\lambda(p))\le\kappa_b$ for
$\gg b^{\lambda}\lambda^{-2}$ primes $p\in\mathcal P_\lambda$, with $\kappa_{10}=1378$. It is not known
whether $R_\lambda(p)$ is prime for infinitely many primes $p$, nor whether there are infinitely many
palindromic primes.

\subsection{Notation}\label{sec:notation}

The letters $p$ and $\ell$ always denote primes, $b\ge2$ is fixed, and $\lambda\ge2$ is an integer
parameter, which is understood throughout to tend to infinity. We set $\omega(1)=\Omega(1)=0$, write
$v_\ell(n)$ for the exponent of $\ell$ in the factorisation of $n$, and put
\[
\Phi(t)\ :=\ \frac1{\sqrt{2\pi}}\int_{-\infty}^{t}e^{-u^{2}/2}\,du .
\]
Implied constants in the Landau and Vinogradov symbols $O(\cdot)$, $o(\cdot)$, $\ll$ and $\asymp$ may
depend on $b$ and on the level-of-distribution exponent $\xi$ of Proposition~\ref{prop:criterion}, and
in Section~\ref{sec:criterion} also on $\delta$, $\#\mathcal E$ and the implied constant in
\eqref{eq:regular} (all introduced in \S\ref{sec:results} below); they depend on other parameters only
where these are displayed as subscripts. In particular no implied constant depends on $k$.

\subsection{Statement of results}\label{sec:results}

Throughout we abbreviate
\begin{equation}\label{eq:defCk}
L:=\log\log b^{\lambda},\qquad C_k:=\frac{\Gamma(k+1)}{2^{k/2}\,\Gamma(k/2+1)},
\end{equation}
so that $C_k$ is the $k$-th moment of the standard normal law when $k$ is even. Every $n$ occurring in
Theorems~\ref{thm:pal}--\ref{thm:moments} lies in $[b^{\lambda-1},b^{\lambda})$ for $\lambda\ge3$, as
the proof of Corollary~\ref{cor:normalorder} records, so that $L=\log\log n+O(\lambda^{-1})$; the
reader may therefore read $L$ as $\log\log n$ throughout.

\begin{theorem}\label{thm:pal}
Let $b\ge2$ be an integer. Then, uniformly for $t\in\mathbb R$,
\[
\frac1{\#\mathcal T_\lambda}\#\Bigl\{n\in\mathcal T_\lambda:\ \omega(n)\le L+t\sqrt L\Bigr\}
\ =\ \Phi(t)+o(1)\qquad(\lambda\to\infty),
\]
and the same asymptotic holds with $\omega$ replaced by $\Omega$.
\end{theorem}


For $1\le i\le b-1$ put
\[
\mathcal P_{\lambda,i}:=\bigl\{p\in\mathcal P_\lambda:\ \eps_{\lambda-1}(p)=i\bigr\}
=\bigl\{p\ \text{prime}:\ ib^{\lambda-1}\le p<(i+1)b^{\lambda-1}\bigr\},
\qquad \pi_{\lambda,i}:=\#\mathcal P_{\lambda,i},
\]
so that $\mathcal P_\lambda$ is the disjoint union of the $\mathcal P_{\lambda,i}$, and write
\begin{equation}\label{eq:defA}
\mathcal A_\lambda:=\bigl\{R_\lambda(p):\ p\in\mathcal P_\lambda\bigr\},\qquad
\mathcal A_{\lambda,i}:=\bigl\{R_\lambda(p):\ p\in\mathcal P_{\lambda,i}\bigr\}.
\end{equation}
By \eqref{eq:defR}, $i$ is the last digit of $R_\lambda(p)$ for $p\in\mathcal P_{\lambda,i}$.
Reversal of a block of $\lambda$ digits is an involution, so that $R_\lambda$ is injective on
$\mathcal P_\lambda$ (see the start of Section~\ref{sec:local}); hence
$\#\mathcal A_\lambda=\pi_\lambda$ and $\#\mathcal A_{\lambda,i}=\pi_{\lambda,i}$, and averages over
$\mathcal A_\lambda$ coincide with averages over $\mathcal P_\lambda$.

\begin{theorem}\label{thm:main}
Let $b\ge2$ be an integer and let $f$ denote either $\omega$ or $\Omega$. Then, uniformly for
$t\in\mathbb R$,
\[
\frac{1}{\pi_\lambda}\,
\#\Bigl\{p\in\mathcal P_\lambda:\ f\bigl(R_\lambda(p)\bigr)\le L+t\sqrt L\Bigr\}\ =\ \Phi(t)+o(1)
\qquad(\lambda\to\infty),
\]
and the same asymptotic holds for each fixed $i\in\{1,\dots,b-1\}$ with $\mathcal P_\lambda$ and
$\pi_\lambda$ replaced throughout by $\mathcal P_{\lambda,i}$ and $\pi_{\lambda,i}$.
\end{theorem}

The localised form of Theorem~\ref{thm:main} is the Erd\H os--Kac law for the reversals of the primes
whose leading digit is prescribed. The unlocalised form follows by summing over $i$.

The statements for $\omega$ in Theorems~\ref{thm:pal} and \ref{thm:main} both follow from a moment
estimate that is uniform in the order of the moment.

\begin{theorem}\label{thm:moments}
Let $b\ge2$ and let $\mathcal B_\lambda$ be $\mathcal T_\lambda$, or $\mathcal A_\lambda$, or
$\mathcal A_{\lambda,i}$ for a fixed $i\in\{1,\dots,b-1\}$. Then, uniformly for positive integers
$k\le\tfrac12L^{1/3}$, one has
\[
\begin{aligned}
\frac1{\#\mathcal B_\lambda}\sum_{n\in\mathcal B_\lambda}\bigl(\omega(n)-L\bigr)^{k}
  &\ =\ C_k\,L^{k/2}\Bigl(1+O\bigl(k^{3/2}L^{-1/2}\bigr)\Bigr) && (k\ \text{even}),\\
\frac1{\#\mathcal B_\lambda}\sum_{n\in\mathcal B_\lambda}\bigl(\omega(n)-L\bigr)^{k}
  &\ \ll\ C_kL^{k/2}k^{3/2}L^{-1/2}                            && (k\ \text{odd}).
\end{aligned}
\]
\end{theorem}

For $\mathcal B_\lambda$ the set of all integers below $b^{\lambda}$, Theorem~\ref{thm:moments} is
contained in Theorem 1 of Granville and Soundararajan \cite{GS}, from which the form of the error term
is taken. The exponent $1/3$ is the one imposed by their moment estimate (Proposition~\ref{prop:gs}
below); the explicit factor $\frac12$ is a convenience, absorbing the difference between $L$ and the
variance of the sieve model constructed in Section~\ref{sec:criterion}. Theorems~\ref{thm:pal} and
\ref{thm:main} use only the case of fixed $k$; the uniformity is recorded with a view to
Remark~\ref{rem:rate}. Theorem~\ref{thm:moments} is stated for $\omega$ only: the corresponding
statements for $\Omega$ in Theorems~\ref{thm:pal} and \ref{thm:main} come from hypothesis (iii) of
Proposition~\ref{prop:criterion} below and are not quantitative.

\begin{corollary}\label{cor:normalorder}
For every $\theta>0$ the number of $n\in\mathcal T_\lambda$ with $|\omega(n)-L|>\theta L$ is
$o(\#\mathcal T_\lambda)$, and the number of $p\in\mathcal P_\lambda$ with
$|\omega(R_\lambda(p))-L|>\theta L$ is $o(\pi_\lambda)$; likewise for $\Omega$. In particular,
$\omega$ and $\Omega$ have normal order $\log\log n$ on both families.
\end{corollary}

In particular, for each fixed $K$, all but $o(\pi_\lambda)$ primes $p\in\mathcal P_\lambda$ have
$\omega(R_\lambda(p))>K$; so the set of $p\in\mathcal P_\lambda$ with $\Omega(R_\lambda(p))\le\kappa_b$
has density zero in $\mathcal P_\lambda$. This is consistent with the lower bound
$\gg b^{\lambda}\lambda^{-2}$ for the cardinality of that set obtained in \cite{DRS}, since
$\pi_\lambda\asymp b^{\lambda}\lambda^{-1}$.
The constant $\kappa_b$ invites comparison with the normal order, though the comparison is one of
sizes only, Corollary~\ref{cor:normalorder} being asymptotic and carrying no information for a
bounded range of $\lambda$. (For $b=10$ one has $L>\kappa_{10}=1378$ only when
$\lambda>e^{1378}/\log10$.) What makes the count of \cite{DRS} a statement about reversals with
atypically few prime factors is that it holds for every large $\lambda$ with $\kappa_b$ held fixed.

Theorems~\ref{thm:pal}, \ref{thm:main} and \ref{thm:moments} are all deduced from the following
criterion. Let $(\mathcal B_\lambda)_{\lambda\ge2}$ be
non-empty finite sets of positive integers, put $N_\lambda:=\#\mathcal B_\lambda$, and for $d\ge1$ set
\begin{equation}\label{eq:defr}
r_d(\lambda)\ :=\ \#\bigl\{n\in\mathcal B_\lambda:\ d\mid n\bigr\}-\frac{N_\lambda}{d}.
\end{equation}
For a finite set of primes $\mathcal E$ let
\begin{equation}\label{eq:defD}
\mathcal D(\mathcal E)\ :=\ \bigl\{d\ge1:\ \ell\mid d\ \Rightarrow\ \ell\notin\mathcal E\bigr\}.
\end{equation}
Finally, call a set $S$ of primes \emph{regular of density $\delta$}, where $0<\delta\le1$, if
\begin{equation}\label{eq:regular}
\sum_{\substack{\ell\le y\\ \ell\in S}}\frac1\ell\ =\ \delta\log\log y+O(1)\qquad(y\ge3),
\end{equation}
and put
\[
\omega_S(n):=\#\{\ell\in S:\ \ell\mid n\},\qquad \Omega_S(n):=\sum_{\ell\in S}v_\ell(n),
\]
so that $\omega_S=\omega$ and $\Omega_S=\Omega$ when $S$ is the set of all primes, which is regular of
density $1$ by Mertens' theorem.

\begin{proposition}[cf.\ {\cite[Prop.~3]{GS}}, {\cite[Thm.~3]{MMP}}]\label{prop:criterion}
Let $(\mathcal B_\lambda)_{\lambda\ge2}$, $N_\lambda$, $r_d(\lambda)$ and $\mathcal D(\mathcal E)$ be as
above. Suppose that there exist a constant $\xi>0$ and a finite set of primes $\mathcal E$, both
independent of $\lambda$, such that
\begin{enumerate}
\item[\textup{(i)}] $n\le b^{\lambda}$ for every $n\in\mathcal B_\lambda$, and
\item[\textup{(ii)}] for every $A>0$,
\[
\sum_{\substack{d\le b^{\xi\lambda},\ d\ \text{squarefree}\\ d\in\mathcal D(\mathcal E)}}
\bigl|r_d(\lambda)\bigr|\ \ll_A\ N_\lambda\,\lambda^{-A}.
\]
\end{enumerate}
Let $S$ be regular of density $\delta$. Then, uniformly for positive integers
$k\le\frac12(\delta L)^{1/3}$, one has
\begin{equation}\label{eq:momk}
\begin{aligned}
\frac1{N_\lambda}\sum_{n\in\mathcal B_\lambda}\bigl(\omega_S(n)-\delta L\bigr)^{k}
  &\ =\ C_k(\delta L)^{k/2}\Bigl(1+O\bigl(k^{3/2}(\delta L)^{-1/2}\bigr)\Bigr) && (k\ \text{even}),\\
\frac1{N_\lambda}\sum_{n\in\mathcal B_\lambda}\bigl(\omega_S(n)-\delta L\bigr)^{k}
  &\ \ll\ C_k(\delta L)^{k/2}k^{3/2}(\delta L)^{-1/2}                          && (k\ \text{odd}).
\end{aligned}
\end{equation}
Consequently, uniformly for $t\in\mathbb R$,
\begin{equation}\label{eq:concl}
\frac1{N_\lambda}\#\Bigl\{n\in\mathcal B_\lambda:\ \omega_S(n)\le\delta L
+t\sqrt{\delta L}\Bigr\}\ =\ \Phi(t)+o(1)\qquad(\lambda\to\infty).
\end{equation}
If in addition
\begin{enumerate}
\item[\textup{(iii)}] $\displaystyle\sum_{n\in\mathcal B_\lambda}\bigl(\Omega(n)-\omega(n)\bigr)\ \ll\ N_\lambda$,
\end{enumerate}
then \eqref{eq:concl} also holds with $\omega_S$ replaced by $\Omega_S$.
\end{proposition}

The constant $\xi$ is any exponent of distribution available for the family: for palindromes any
$\xi<\beta(b)$ of Theorem~\ref{thm:col}, for reversed primes any $\xi<\xi_0(b)$ of
Theorem~\ref{thm:drs}. The conclusion does not depend on its value, only on its existence
(Remark~\ref{rem:inputs}). Theorems~\ref{thm:pal}--\ref{thm:moments} are the case in which $S$ is the
set of all primes and $\delta=1$; Corollary~\ref{cor:cheb} below is the general case.

Proposition~\ref{prop:criterion} differs from \cite[Prop.~3]{GS} and \cite[Thm.~3]{MMP} in exactly one
respect: the presence of the finite set $\mathcal E$, at whose primes no hypothesis is imposed. That
single change is what makes the two families of this note accessible. In a typical application of the
moment estimate of \cite{GS} (quoted as Proposition~\ref{prop:gs} below) one chooses the multiplicative
function $h$ appearing there so as to model the local densities of the family; the hypothesis
$0\le h(d)\le d$ then encodes what is known about them. In the proof of
Proposition~\ref{prop:criterion} nothing is modelled. We take $h\equiv1$, so that $r_d$ means exactly
\eqref{eq:defr}, and hypothesis (ii) alone carries the assertion that the remainders are small on
average. Both families treated here need this latitude: by Lemma~\ref{lem:dens}(i) and
Example~\ref{ex:base10} the local densities of $R_\lambda(p)$ at the primes dividing $b$ are not
$1/\ell$, and by Lemma~\ref{lem:cong}(i) at those dividing $b^{2}-1$ they vanish, while for palindromes
the constraints are those described after \eqref{eq:defE}; any criterion that insists on modelling them
correctly is unusable here. Since $\mathcal E$ is finite and independent of $\lambda$, discarding its
primes shifts $\omega_S$ by $O(1)$, which is $o(\sqrt L)$ and therefore invisible on the scale of
\eqref{eq:concl}: whatever happens at those primes cannot affect the limit law.

The proof of Proposition~\ref{prop:criterion} runs as follows. Given $k$, one truncates at
$y=b^{\xi\lambda/(k+1)}$, so that products of $k+1$ primes below $y$ stay within the level
$b^{\xi\lambda}$ of hypothesis (ii); the truncation is cheap because
$\log\log y=L+\log\xi-\log(k+1)$, so that it moves the mean by $O(\log(2k))$ only. The primes of
$\mathcal E$ and those exceeding $y$ are then discarded, which by hypothesis (i) alters $\omega_S(n)$ by
at most $\#\mathcal E+\lambda\log b/\log y\ll k=o(\sqrt L)$. What is left is a sum over primes at which
hypothesis (ii) is available, and to it one applies the estimate of Granville and Soundararajan.

Taking for $S$ the primes with a prescribed Frobenius class gives the following.

\begin{corollary}\label{cor:cheb}
Let $K$ be a finite Galois extension of $\mathbb Q$, let $\mathcal C$ be a conjugacy class in
$\Gal(K/\mathbb Q)$, and let $S_{\mathcal C}$ be the set of primes unramified in $K$ whose Frobenius
class is $\mathcal C$; put $\delta:=\#\mathcal C/[K:\mathbb Q]$. Then Theorems~\ref{thm:pal},
\ref{thm:main} and \ref{thm:moments} hold with $\omega$ replaced by $\omega_{S_{\mathcal C}}$ and
$\Omega$ by $\Omega_{S_{\mathcal C}}$, and with $L$ and $\sqrt L$ replaced by $\delta L$ and
$\sqrt{\delta L}$; in particular the moments are uniform for positive integers
$k\le\frac12(\delta L)^{1/3}$, with error term $O\bigl(k^{3/2}(\delta L)^{-1/2}\bigr)$.
\end{corollary}

Proposition~\ref{prop:criterion} is proved in Section~\ref{sec:criterion}. Section~\ref{sec:pal} treats
palindromes and Section~\ref{sec:local} reversed primes; the latter uses the set $\mathcal E_b$ of
\eqref{eq:defE} and one identity from the proof of Lemma~\ref{lem:palhyp}, but is otherwise independent
of the former. Section~\ref{sec:remarks} collects four remarks: on the arithmetic inputs and why a
Siegel--Walfisz range would not suffice (Remark~\ref{rem:inputs}), on the extremal problem of
\cite{BSh} (Remark~\ref{rem:extremal}), on other levels of distribution for palindromes
(Remark~\ref{rem:paldist}), and on rates of convergence (Remark~\ref{rem:rate}).


\section{An Erd\H os--Kac criterion}\label{sec:criterion}

We use the sieve-theoretic moment estimate of Granville and Soundararajan in the following form. Let
$\mathcal M$ be a finite multiset of $N$ positive integers (we shall use it only when $\mathcal M$ is a
set), put $N(d):=\#\{n\in\mathcal M:d\mid n\}$, let $h$ be a non-negative multiplicative function with
$0\le h(d)\le d$ for squarefree $d$, and write $N(d)=h(d)N/d+r_d$ for squarefree $d$. For a set of
primes $\mathcal R$ put
\[
\mu_{\mathcal R}:=\sum_{\ell\in\mathcal R}\frac{h(\ell)}{\ell},\qquad
\sigma_{\mathcal R}^{2}:=\sum_{\ell\in\mathcal R}\frac{h(\ell)}{\ell}\Bigl(1-\frac{h(\ell)}{\ell}\Bigr),
\]
these being the mean and the variance of the sieve model introduced after \eqref{eq:defy} below; put
$\omega_{\mathcal R}(n):=\#\{\ell\in\mathcal R:\ \ell\mid n\}$, as before; and, for an integer
$j\ge0$, let $\Pi_j(\mathcal R)$ be the set of squarefree integers which are products of at most $j$
primes of $\mathcal R$. The letter $\mu$ is used for this mean throughout, the M\"obius function not
occurring in this note.

\begin{proposition}[Granville and Soundararajan {\cite[Prop.~3]{GS}}]\label{prop:gs}
Uniformly for all positive integers $k\le\sigma_{\mathcal R}^{2/3}$, one has
\[
\begin{aligned}
\sum_{n\in\mathcal M}\bigl(\omega_{\mathcal R}(n)-\mu_{\mathcal R}\bigr)^{k}
  &= C_kN\sigma_{\mathcal R}^{k}\Bigl(1+O\Bigl(\frac{k^{3}}{\sigma_{\mathcal R}^{2}}\Bigr)\Bigr)
     +O\Bigl(\mu_{\mathcal R}^{k}\!\!\sum_{d\in \Pi_k(\mathcal R)}\!\!|r_d|\Bigr)
  && (k\ \text{even}),\\
\sum_{n\in\mathcal M}\bigl(\omega_{\mathcal R}(n)-\mu_{\mathcal R}\bigr)^{k}
  &\ll C_kN\sigma_{\mathcal R}^{k}\frac{k^{3/2}}{\sigma_{\mathcal R}}
     +\mu_{\mathcal R}^{k}\!\!\sum_{d\in \Pi_k(\mathcal R)}\!\!|r_d|
  && (k\ \text{odd}).
\end{aligned}
\]
\end{proposition}

Throughout this section $(\mathcal B_\lambda)$, $N_\lambda$, $\xi$, $\mathcal E$ and $S$ are as
in Proposition~\ref{prop:criterion}, with $S$ regular of density $\delta$. Given a positive integer
$k\le\frac12(\delta L)^{1/3}$, set
\begin{equation}\label{eq:defy}
y\ :=\ b^{\xi\lambda/(k+1)},\qquad
\mathcal Q\ :=\ \{\ell\le y:\ \ell\in S,\ \ell\notin\mathcal E\},\qquad
\mu:=\mu_{\mathcal Q},\quad \sigma:=\sigma_{\mathcal Q},
\end{equation}
the last two formed with $h\equiv1$, so that $\mu=\sum_{\ell\in\mathcal Q}\ell^{-1}$ and
$\sigma^{2}=\sum_{\ell\in\mathcal Q}\ell^{-1}(1-\ell^{-1})$. These are the mean and the variance of
$\sum_{\ell\in\mathcal Q}X_\ell$, where the $X_\ell$ are independent Bernoulli variables with
$\mathbb P(X_\ell=1)=1/\ell$: this is the sieve model for $\omega_{\mathcal Q}$ in the sense of
\cite{GS}, and Proposition~\ref{prop:gs} says that the moments of $\omega_{\mathcal Q}-\mu$ follow it.

The next lemma records the three properties of \eqref{eq:defy} that the proof needs: that the
truncation moves the mean by $O(\log(2k))$ only; that discarding $\mathcal E$ and the primes above $y$
changes $\omega_S$ by $O(k)$; and that every modulus arising in Proposition~\ref{prop:gs} is covered by
hypothesis (ii).

\begin{lemma}\label{lem:setup}
With the notation \eqref{eq:defy} one has, uniformly for $1\le k\le\frac12(\delta L)^{1/3}$:
\begin{enumerate}
\item[\textup{(i)}] $\mu=\delta L+O(\log(2k))$ and $\sigma^{2}=\delta L+O(\log(2k))$;
\item[\textup{(ii)}] $\bigl|\omega_S(n)-\omega_{\mathcal Q}(n)\bigr|\le (k+1)/\xi+\#\mathcal E\ \ll\ k$ for every
$n\in\mathcal B_\lambda$;
\item[\textup{(iii)}] $\displaystyle\sum_{d\in \Pi_{k+1}(\mathcal Q)}|r_d(\lambda)|\ \ll_A\ N_\lambda\lambda^{-A}$
for every $A>0$.
\end{enumerate}
\end{lemma}

\begin{proof}
(i) By \eqref{eq:regular}, $\sum_{\ell\le y,\ \ell\in S}\ell^{-1}=\delta\log\log y+O(1)$, and deleting the
at most $\#\mathcal E$ primes of $\mathcal E$ alters this by $O(1)$; hence $\mu=\delta\log\log y+O(1)$.
From $\log y=\bigl(\xi/(k+1)\bigr)\lambda\log b$ we get $\log\log y=L+\log\xi-\log(k+1)$, which
gives the first assertion. The second follows from
$\sigma^{2}=\mu-\sum_{\ell\in\mathcal Q}\ell^{-2}=\mu+O(1)$.

(ii) The difference is at most $\#\{\ell\in\mathcal E:\ell\mid n\}+\#\{\ell>y:\ell\mid n\}$. The first term
is at most $\#\mathcal E$; for the second, $n\le b^{\lambda}$ by hypothesis (i) and every prime counted
exceeds $y=b^{\xi\lambda/(k+1)}$, so their number is less than $\lambda\log b/\log y=(k+1)/\xi$.

(iii) Recall that $\Pi_{k+1}(\mathcal Q)$ is the set of squarefree integers which are products of at
most $k+1$ primes of $\mathcal Q$, and let $d=\ell_1\cdots\ell_j\in \Pi_{k+1}(\mathcal Q)$, the
$\ell_i$ being distinct elements of $\mathcal Q$ and $j\le k+1$. Then $d$ is squarefree; every prime
factor of $d$ lies in $\mathcal Q$ and hence outside $\mathcal E$, so that
$d\in\mathcal D(\mathcal E)$ by \eqref{eq:defD}; and each $\ell_i\le y$, so that
\[
d\ \le\ y^{j}\ \le\ y^{k+1}\ =\ b^{\xi\lambda}
\]
by the choice of $y$ in \eqref{eq:defy}. Hence
\[
\Pi_{k+1}(\mathcal Q)\ \subseteq\
\bigl\{d\le b^{\xi\lambda}:\ d\ \text{squarefree},\ d\in\mathcal D(\mathcal E)\bigr\},
\]
and since the terms $|r_d(\lambda)|$ are non-negative, the sum in (iii) is at most the sum in
hypothesis (ii) of Proposition~\ref{prop:criterion}, which is $\ll_A N_\lambda\lambda^{-A}$. This is
what dictates the exponent in \eqref{eq:defy}: the truncation point $y$ is chosen so that $y^{k+1}$
is exactly the level $b^{\xi\lambda}$ supplied by hypothesis (ii), with $k+1$ rather than $k$
because Proposition~\ref{prop:gs} is applied below at every order $m\le k+1$, and
$\Pi_m(\mathcal Q)\subseteq \Pi_{k+1}(\mathcal Q)$ for such $m$.
\end{proof}

We can now say where the range $k\le\frac12(\delta L)^{1/3}$ of Proposition~\ref{prop:criterion} comes
from. It arises twice. Firstly, Proposition~\ref{prop:gs} is available only for
$k\le\sigma^{2/3}$, and by Lemma~\ref{lem:setup}(i) the variance of the sieve model satisfies
$\sigma^{2}=\delta L+O(\log(2k))$. Secondly, the binomial estimate in the proof below requires
$k^{3/2}\ll\sigma$, which is the same condition. Both therefore read $k\ll(\delta L)^{1/3}$, the factor
$\frac12$ absorbing the $O(\log(2k))$ discrepancy between $\sigma^{2}$ and $\delta L$, together with the
passage from $k$ to $k+1$ in the proof.

\begin{proof}[Proof of Proposition~\ref{prop:criterion}]
Fix $k\le\frac12(\delta L)^{1/3}$ and let $y,\mathcal Q,\mu,\sigma$ be as in \eqref{eq:defy}. Apply
Proposition~\ref{prop:gs} with $\mathcal M=\mathcal B_\lambda$, $N=N_\lambda$, $h\equiv1$ and
$\mathcal R=\mathcal Q$; the remainders $r_d$ are then those of \eqref{eq:defr}, and only
$d\in\mathcal D(\mathcal E)$ occur in the conclusion, every element of $\Pi_{k+1}(\mathcal Q)$ being a
product of primes of $\mathcal Q$ and hence coprime to $\mathcal E$. By Lemma~\ref{lem:setup}(i) we have
$\sigma^{2/3}=(\delta L)^{1/3}(1+o(1))\ge k+1$ for all large $\lambda$, uniformly for $k$ in the
stated range, so Proposition~\ref{prop:gs} is available for every order $m\le k+1$.

Write $W(n):=\omega_{\mathcal Q}(n)-\mu$. By Lemma~\ref{lem:setup}(i) there is a constant $C$, not
depending on $k$, with $\mu\le\delta L+C\log(2k)\le L\bigl(1+C\log(2k)/L\bigr)$, since $\delta\le1$.
Raising this to the $m$-th power and using $1+u\le e^{u}$ gives
$\mu^{m}\le L^{m}\exp\bigl(Cm\log(2k)/L\bigr)$; and for $m\le k+1$ the exponent is $o(1)$, because
$2k\le(\delta L)^{1/3}$ gives $\log(2k)\le\frac13\log L$ while $m\ll L^{1/3}$, so that
$m\log(2k)\ll L^{1/3}\log L=o(L)$. Thus $\mu^{m}\le L^{m}e^{o(1)}$. In particular no constant is
raised here to a power that grows with $k$. With Lemma~\ref{lem:setup}(iii) at $A=1$,
the second error term of Proposition~\ref{prop:gs}
is therefore $\ll L^{m}\lambda^{-1}$ after division by $N_\lambda$. Since $\lambda=e^{L}/\log b$ and
$m\le k+1\ll L^{1/3}$, this is $e^{-L(1+o(1))}$, hence smaller than any fixed power of
$L^{-1}$. As $C_m\sigma^{m}\ge1$ for $\lambda$ large, and since $\sigma^{2}\asymp L$, it is in
particular smaller than both $C_m\sigma^{m}m^{3}\sigma^{-2}$ and $C_m\sigma^{m}m^{3/2}\sigma^{-1}$, and so
may be absorbed into the error terms below. Hence, for $1\le m\le k+1$,
\begin{equation}\label{eq:D0}
\begin{aligned}
\frac1{N_\lambda}\sum_{n}W(n)^{m}
  &=   C_m\sigma^{m}\Bigl(1+O\bigl(m^{3}\sigma^{-2}\bigr)\Bigr) && (m\ \text{even}),\\
\frac1{N_\lambda}\sum_{n}W(n)^{m}
  &\ll C_m\sigma^{m}m^{3/2}\sigma^{-1}                          && (m\ \text{odd}).
\end{aligned}
\end{equation}
In particular, for every $0\le m\le k$,
\begin{equation}\label{eq:absmom}
\frac1{N_\lambda}\sum_{n}\bigl|W(n)\bigr|^{m}\ \ll\ C_m\sigma^{m}.
\end{equation}
For $m=0$ both sides equal $1$; for $1\le m\le k$, Lyapunov's inequality and the even case of
\eqref{eq:D0} give
\[
\frac1{N_\lambda}\sum_{n}\bigl|W(n)\bigr|^{m}
\ \le\ \Bigl(\frac1{N_\lambda}\sum_{n}W(n)^{2\lceil m/2\rceil}\Bigr)^{m/(2\lceil m/2\rceil)}
\ \ll\ C_m\sigma^{m},
\]
where for odd $m$ the last step uses
$C_{m+1}^{m/(m+1)}\ll C_m$, which follows from $C_m=\Gamma(m+1)/(2^{m/2}\Gamma(m/2+1))$ and Stirling's
formula.

Now put $\Delta(n):=\bigl(\omega_S(n)-\omega_{\mathcal Q}(n)\bigr)+\bigl(\mu-\delta L\bigr)$, so that
$\omega_S(n)-\delta L=W(n)+\Delta(n)$ and, by Lemma~\ref{lem:setup}(i) and (ii),
$\|\Delta\|_\infty\ll k+\log(2k)\ll k$. By the binomial theorem,
\[
\frac1{N_\lambda}\sum_{n}\bigl(\omega_S(n)-\delta L\bigr)^{k}
=\frac1{N_\lambda}\sum_{n}W(n)^{k}
+\sum_{j=1}^{k}\binom kj\frac1{N_\lambda}\sum_{n}W(n)^{k-j}\Delta(n)^{j},
\]
and by \eqref{eq:absmom} the $j$-th summand is $\ll\binom kj(c_1k)^{j}C_{k-j}\sigma^{k-j}$, where $c_1$
depends only on the parameters held fixed in this section and, in particular, not on $k$. Stirling's
formula gives $C_m\asymp(m/e)^{m/2}$ for $m\ge1$, and $C_0=1$; hence
$C_{k-j}\ll C_k(e/k)^{j/2}$ for $0\le j\le k$, because
$\bigl((k-j)/k\bigr)^{(k-j)/2}\le1$. So, using $\binom kj\le k^{j}/j!$ and writing
$c_2:=c_1\sqrt e$ and absorbing the implied constant just introduced,
\[
\sum_{j=1}^{k}\binom kj(c_1k)^{j}C_{k-j}\sigma^{k-j}
\ \ll\ C_k\sigma^{k}\sum_{j\ge1}\frac1{j!}\Bigl(\frac{c_2k^{3/2}}{\sigma}\Bigr)^{j}
\ \ll\ C_k\sigma^{k}\,\frac{k^{3/2}}{\sigma},
\]
the last step because $u:=c_2k^{3/2}/\sigma=O(1)$ in the stated range, whence
$e^{u}-1\ll u$. Finally, write
$\sigma^{2}=\delta L(1+\vartheta)$, where $\vartheta\ll\log(2k)/(\delta L)$ by
Lemma~\ref{lem:setup}(i). Then $k|\vartheta|\ll(\delta L)^{-2/3}\log(\delta L)=o(1)$ uniformly in
the stated range, so
\[
\sigma^{k}=(\delta L)^{k/2}\exp\bigl(\tfrac k2\log(1+\vartheta)\bigr)
=(\delta L)^{k/2}\bigl(1+O(k\log(2k)/(\delta L))\bigr),
\]
and $k\log(2k)/(\delta L)\ll k^{3/2}(\delta L)^{-1/2}$. Combining this with \eqref{eq:D0} for $m=k$,
and noting $k^{3}\sigma^{-2}\ll k^{3/2}\sigma^{-1}$ in the stated range, gives \eqref{eq:momk} and its odd
counterpart.

For \eqref{eq:concl}, fix $k$ and let $\lambda\to\infty$ in \eqref{eq:momk}: the $k$-th moment of
$(\omega_S-\delta L)/\sqrt{\delta L}$ tends to $C_k$ for $k$ even and to $0$ for $k$ odd, which
are the moments of the standard normal law. That law is determined by its moments, so \eqref{eq:concl}
follows by the method of moments in the form of \cite[Thm.~30.2]{Bil}, uniformly in $t$ because $\Phi$ is
continuous.

Assume now hypothesis (iii), and let $F^{\omega}_\lambda$ and $F^{\Omega}_\lambda$ denote the
distribution functions in \eqref{eq:concl} formed with $\omega_S$ and with $\Omega_S$ respectively.
Since $\Omega_S(n)-\omega_S(n)=\sum_{\ell\in S}\max(v_\ell(n)-1,0)\le\Omega(n)-\omega(n)$, hypothesis
(iii) gives $N_\lambda^{-1}\sum_n(\Omega_S(n)-\omega_S(n))\ll1$. Fix $\tau>0$ and let
$\eta_\tau(\lambda)$ be the proportion of $n\in\mathcal B_\lambda$ with
$\Omega_S(n)-\omega_S(n)>\tau\sqrt{\delta L}$; by Markov's inequality applied to the non-negative
function $\Omega_S-\omega_S$ we have $\eta_\tau(\lambda)\ll1/(\tau\sqrt{\delta L})$, a bound depending
on $\lambda$ and $\tau$ but not on $t$. Now $\Omega_S\ge\omega_S$ pointwise, so for every real $c$ the
condition $\Omega_S(n)\le c$ forces $\omega_S(n)\le c$; and conversely
$\omega_S(n)\le c-\tau\sqrt{\delta L}$ forces $\Omega_S(n)\le c$ unless $n$ is one of the
$\eta_\tau(\lambda)N_\lambda$ exceptions. Taking $c=\delta L+t\sqrt{\delta L}$ therefore gives
\[
F^{\omega}_\lambda(t-\tau)-\eta_\tau(\lambda)\ \le\ F^{\Omega}_\lambda(t)
\ \le\ F^{\omega}_\lambda(t)\qquad(t\in\mathbb R),
\]
whence
\[
\sup_t\bigl|F^{\Omega}_\lambda(t)-\Phi(t)\bigr|\ \le\ \sup_t\bigl|F^{\omega}_\lambda(t)-\Phi(t)\bigr|
+\sup_t\bigl|\Phi(t)-\Phi(t-\tau)\bigr|+\eta_\tau(\lambda),
\]
in which the middle term is at most $\tau/\sqrt{2\pi}$, since $\Phi'\le(2\pi)^{-1/2}$. Letting
$\lambda\to\infty$ with $\tau$ fixed, and then $\tau\to0$, gives \eqref{eq:concl} for $\Omega_S$,
uniformly in $t$.
\end{proof}

\section{Palindromes}\label{sec:pal}

Write $\mathcal T:=\bigcup_{\nu\ge1}\mathcal T_\nu$ for the set of all base-$b$ palindromes and, for
$x\ge1$, $a\in\mathbb Z$ and $q\ge1$, put $\mathcal T(x):=\{n\in\mathcal T:n<x\}$ and
$\mathcal T(x,a,q):=\{n\in\mathcal T(x):n\equiv a\ (\mathrm{mod}\ q)\}$. Since
$\#\mathcal T_\nu=(b-1)b^{\lceil\nu/2\rceil-1}$, summation over $\nu\le\lambda$ gives
\begin{equation}\label{eq:Tsize}
\#\mathcal T\bigl(b^{\lambda}\bigr)\ \asymp\ b^{\lceil\lambda/2\rceil}\ \asymp\ \#\mathcal T_\lambda ,
\end{equation}
since the family grows geometrically, so that the cumulative and the per-length counts are
interchangeable; this is what makes the passage from $\mathcal T(x)$ to $\mathcal T_\lambda$ in
Lemma~\ref{lem:palhyp} harmless. We use two results from the literature: a level of distribution, and
an upper bound of Brun--Titchmarsh type valid for every modulus.

\begin{theorem}[Col {\cite[Thm.~2]{Col}}]\label{thm:col}
For every integer $b\ge2$ there is a constant $\beta=\beta(b)>0$ such that for all $A>0$ and $\eta>0$,
\[
\sum_{\substack{q<x^{\beta-\eta}\\ \gcd(q,\,b^{3}-b)=1}}\ \sup_{z\le x}\ \max_{a\in\mathbb Z}
\left|\ \#\mathcal T(z,a,q)-\frac{\#\mathcal T(z)}{q}\ \right|
\ \ll_{b,A,\eta}\ \frac{\#\mathcal T(x)}{\log^{A}x}.
\]
\end{theorem}

Admissible values of $\beta(b)$ are tabulated in \cite[\S1]{Col}: one may take $\beta(2)=1/30$ and
$\beta(10)=1/186$, and $\beta(b)\sim1/(6\pi b)$ as $b\to\infty$.

\begin{theorem}[Banks and Shparlinski {\cite[Thm.~7]{BSh}}]\label{thm:bsh}
For every integer $b\ge2$, every $\nu\ge1$ and every integer $d\ge1$,
$\ \#\{n\in\mathcal T_\nu:\ d\mid n\}\ \ll_b\ \#\mathcal T_\nu\,d^{-1/2}$.
\end{theorem}

Theorem~\ref{thm:bsh} is stated in \cite{BSh} for a fixed number of digits, which is the form used
below. The two results just quoted are complementary:
Theorem~\ref{thm:col} is an equidistribution statement, giving the main term $\#\mathcal T(z)/q$
with a saving of an arbitrary power of $\log x$; but it applies only to $q$ below the level
$x^{\beta-\eta}$ and coprime to $b^{3}-b$, and it supplies hypothesis (ii) of
Proposition~\ref{prop:criterion} directly. Theorem~\ref{thm:bsh} supplies no main term and loses a
square root ($d^{-1/2}$ where equidistribution would give $d^{-1}$), but holds for every $d$,
with neither a coprimality condition nor a restriction of range. It is used only in the verification
of hypothesis (iii), and there only for the moduli beyond the reach of Theorem~\ref{thm:col}. For
these an upper bound suffices, since the quantity being estimated need only be shown to be $O(1)$;
that is what makes the loss of a square root affordable. The exponent $-1/2$ is critical: the case
$m=2$ in the proof of Lemma~\ref{lem:palhyp} produces $\sum_\ell\ell^{-1}$, which is bounded only
because the range is truncated at $b^{\xi\lambda/2}<\ell<b^{\lambda/2}$, the two cut-offs being
comparable powers of $b^{\lambda}$; any weaker exponent would diverge there.

Since $b^{3}-b=b(b^{2}-1)$, the coprimality condition in Theorem~\ref{thm:col} is exactly
$q\in\mathcal D(\mathcal E_b)$ for the set
\begin{equation}\label{eq:defE}
\mathcal E_b\ :=\ \bigl\{\ell:\ \ell\mid b(b^{2}-1)\bigr\},\qquad
\#\mathcal E_b=\omega\bigl(b(b^{2}-1)\bigr)\ll1 .
\end{equation}
The set $\mathcal E_b$ is larger than the structure of $\mathcal T_\lambda$ alone would require. If
$\lambda$ is even then $(b+1)\mid n$ for every $n\in\mathcal T_\lambda$: as $\lambda-1$ is odd, the indices
$j$ and $\lambda-1-j$ are distinct and of opposite parity, so pairing them in
$n\equiv\sum_j\eps_j(n)(-1)^{j}\pmod{b+1}$ gives $n\equiv0$. Moreover the last digit of a palindrome
equals its leading digit and is therefore non-zero, which biases divisibility by the primes dividing
$b$. Theorem~\ref{thm:col} excludes the divisors of $b-1$ as
well. Proposition~\ref{prop:criterion} imposes no hypothesis at the primes of $\mathcal E$, so these
constraints do not enter what follows.

\begin{lemma}\label{lem:palhyp}
Set $\mathcal B_\lambda:=\mathcal T_\lambda$ and $\mathcal E:=\mathcal E_b$, and
let $0<\xi<\beta(b)$. Then hypotheses \textup{(ii)} and \textup{(iii)} of
Proposition~\ref{prop:criterion} hold.
\end{lemma}

\begin{proof}
For (ii), fix $\xi'$ with $\xi<\xi'<\beta(b)$, let $x:=b^{\lambda}$ and let $q\le x^{\xi}$ satisfy
$\gcd(q,b^{3}-b)=1$. Then
\[
\#\bigl\{n\in\mathcal T_\lambda:\ q\mid n\bigr\}=\#\mathcal T(b^{\lambda},0,q)-\#\mathcal T(b^{\lambda-1},0,q),
\qquad
\#\mathcal T_\lambda=\#\mathcal T(b^{\lambda})-\#\mathcal T(b^{\lambda-1}),
\]
both endpoints being at most $x$. Subtracting, the main terms combine to give exactly
$\#\mathcal T_\lambda/q$, so by \eqref{eq:defr}
\[
\bigl|r_q(\lambda)\bigr|\ \le\ 2\sup_{z\le x}\max_{a\in\mathbb Z}
\left|\ \#\mathcal T(z,a,q)-\frac{\#\mathcal T(z)}{q}\ \right| .
\]
As $q\le x^{\xi}<x^{\xi'}$, choosing $\eta:=\beta(b)-\xi'>0$ in Theorem~\ref{thm:col} and summing over $q$
gives $\sum_{q}|r_q(\lambda)|\ll_{b,A}\#\mathcal T(x)(\log x)^{-A}$, which is
$\ll_{b,A}\#\mathcal T_\lambda\lambda^{-A}$ by \eqref{eq:Tsize}. This is (ii), and in fact for all such
$q$ and not only the squarefree ones; we record the stronger form because it will be applied below to
the prime-power moduli $\ell^{m}$ in the verification of hypothesis (iii).

For (iii), use $\Omega(n)-\omega(n)=\sum_{m\ge2}\#\{\ell:\ell^{m}\mid n\}$, valid since a prime $\ell$ is
counted $\max(v_\ell(n)-1,0)$ times on the right; it therefore suffices to show that
\[
\sum_{m\ge2}\ \sum_{\ell}\ \frac{1}{\#\mathcal T_\lambda}
\#\bigl\{n\in\mathcal T_\lambda:\ \ell^{m}\mid n\bigr\}\ \ll\ 1 .
\]
Every prime power occurring divides some $n<b^{\lambda}$,
so $\ell<b^{\lambda/m}$ and $2\le m\le M:=\lambda\log b/\log2$. If $\ell^{m}\le b^{\xi\lambda}$ and
$\ell\nmid b^{3}-b$, then the bound just proved for (ii) applies to the modulus $\ell^{m}$, and after
division by $\#\mathcal T_\lambda$ the total contribution is at
most $\sum_{m\ge2}\sum_{\ell}\ell^{-m}+O(\lambda^{-1})=O(1)$. In every remaining case we appeal to
Theorem~\ref{thm:bsh}, which gives
$\#\{n\in\mathcal T_\lambda:\ell^{m}\mid n\}\ll\#\mathcal T_\lambda\,\ell^{-m/2}$. The primes
$\ell\mid b^{3}-b$ then contribute $\ll\sum_{\ell\mid b^{3}-b}\sum_{m\ge2}\ell^{-m/2}\ll1$, there being
$O(1)$ of them; and for the pairs with $\ell^{m}>b^{\xi\lambda}$ we have $\ell>b^{\xi\lambda/m}$, whence
\[
\begin{aligned}
\sum_{2\le m\le M}\ \sum_{b^{\xi\lambda/m}<\ell<b^{\lambda/m}}\ell^{-m/2}
&\ \ll\ \sum_{b^{\xi\lambda/2}<\ell<b^{\lambda/2}}\frac1\ell
\ +\ \sum_{3\le m\le M}\bigl(b^{\xi\lambda/m}\bigr)^{1-m/2}\\
&\ \ll\ \log\frac1\xi+\lambda\,b^{-\xi\lambda/6}\ \ll\ 1 .
\end{aligned}
\]
Here the first sum is $\log\log b^{\lambda/2}-\log\log b^{\xi\lambda/2}+o(1)=\log(1/\xi)+o(1)$ by Mertens'
theorem, the upper cut-off $\ell<b^{\lambda/2}$ keeping it bounded; in the second sum
$\bigl(b^{\xi\lambda/m}\bigr)^{1-m/2}=b^{-\xi\lambda(m-2)/(2m)}\le b^{-\xi\lambda/6}$ because
$(m-2)/(2m)\ge\frac16$ for $m\ge3$, while the number of terms is $M\ll\lambda$. Adding the three
contributions gives (iii).
\end{proof}

\begin{proof}[Proof of Theorem~\ref{thm:pal}, and of Theorem~\ref{thm:moments} for palindromes]
We apply Proposition~\ref{prop:criterion}, taking $\mathcal B_\lambda:=\mathcal T_\lambda$,
$\mathcal E:=\mathcal E_b$ and $0<\xi<\beta(b)$, with $S$ the set of all primes, so that
$\omega_S=\omega$ and $\delta=1$. Hypothesis (i) holds because every element of $\mathcal T_\lambda$ is
less than $b^{\lambda}$, while hypotheses (ii) and (iii) are Lemma~\ref{lem:palhyp}.
\end{proof}

\section{Reversed primes}\label{sec:local}

We verify the hypotheses of Proposition~\ref{prop:criterion} for the sets $\mathcal A_{\lambda,i}$ of
\eqref{eq:defA}. Reversing a block of $\lambda$ base-$b$ digits, leading zeros retained, is an
involution on the set of such blocks; hence $R_\lambda$ is injective on $\{0,\dots,b^{\lambda}-1\}$,
and in particular on $\mathcal P_\lambda$, so that $\#\mathcal A_{\lambda,i}=\pi_{\lambda,i}$ for
$\lambda\ge2$.

\subsection{Local structure of the reversal}

Reversal is not multiplicative, but it is compatible with reduction modulo $b^{2}-1$ and modulo the
divisors of $b$. The congruences of Lemma~\ref{lem:cong} are those of (3.8) in \cite{DRS}, and its gcd
assertion corresponds to (3.9) there. Stating the first congruence modulo $b^{2}-1$ rather than modulo
$b-1$ and $b+1$ separately is a genuine strengthening when $b$ is odd, since
$\operatorname{lcm}(b-1,b+1)=(b^{2}-1)/2$ then; it also yields the equality of greatest common
divisors.

\begin{lemma}[cf.\ {\cite[(3.8), (3.9)]{DRS}}]\label{lem:cong}
Let $\lambda\ge2$ and let $n$ be an integer with $b^{\lambda-1}\le n<b^{\lambda}$.
\begin{enumerate}
\item[\textup{(i)}] One has $R_\lambda(n)\equiv b^{\lambda-1}n\pmod{b^{2}-1}$; in particular
$R_\lambda(n)\equiv n\pmod{b-1}$ and $R_\lambda(n)\equiv(-1)^{\lambda-1}n\pmod{b+1}$. Moreover
$\gcd\bigl(R_\lambda(n),b^{2}-1\bigr)=\gcd\bigl(n,b^{2}-1\bigr)$, and consequently
$\gcd(R_\lambda(p),b^{2}-1)=1$ for every $p\in\mathcal P_\lambda$ with $p>b+1$. One also has
$b^{\lambda-1}\le R_\lambda(p)<b^{\lambda}$ for every $p\in\mathcal P_\lambda$ with $p>b$.
\item[\textup{(ii)}] If $\ell\mid b$ then $R_\lambda(n)\equiv\eps_{\lambda-1}(n)\pmod\ell$.
\end{enumerate}
\end{lemma}

\begin{proof}
(i) Modulo $b^{2}-1$ we have $b^{2}\equiv1$, hence $b^{-1}\equiv b$ and $b^{-j}\equiv b^{j}$ for every
$j\ge0$; and $\gcd(b,b^{2}-1)=1$, so the powers of $b$ are invertible. Thus \eqref{eq:defR} gives
\[
R_\lambda(n)\ =\ b^{\lambda-1}\sum_{j=0}^{\lambda-1}\eps_j(n)b^{-j}
\ \equiv\ b^{\lambda-1}\sum_{j=0}^{\lambda-1}\eps_j(n)b^{j}
\ =\ b^{\lambda-1}n\pmod{b^{2}-1}.
\]
Reducing modulo $b-1$ and modulo $b+1$ and using $b\equiv1$, respectively $b\equiv-1$, gives the two
stated congruences; and since $b^{\lambda-1}$ is invertible modulo $b^{2}-1$, the displayed congruence
yields $\gcd(R_\lambda(n),b^{2}-1)=\gcd(n,b^{2}-1)$, which is $1$ for a prime $p>b+1$. Finally
$R_\lambda(p)<b^{\lambda}$ is immediate from \eqref{eq:defR}, while the leading digit of $R_\lambda(p)$ is
$\eps_0(p)=p\bmod b$, non-zero because $\gcd(p,b)=1$; hence $R_\lambda(p)\ge b^{\lambda-1}$.

(ii) If $\ell\mid b$ then $b^{\lambda-1-j}\equiv0\pmod\ell$ for $j<\lambda-1$, so \eqref{eq:defR} leaves
only the term with $j=\lambda-1$.
\end{proof}

By \eqref{eq:defD} and \eqref{eq:defE}, $d\in\mathcal D(\mathcal E_b)$ if and only if
$\gcd(d,b(b^{2}-1))=1$. We next compute the local densities of $R_\lambda(p)$ at the primes dividing
$b$; those dividing $b^{2}-1$ were settled in Lemma~\ref{lem:cong}(i).

\begin{lemma}\label{lem:dens}
One has $\pi_{\lambda,i}\sim\pi_\lambda/(b-1)$ for each fixed $i\in\{1,\dots,b-1\}$; in particular
$\pi_{\lambda,i}\asymp\pi_\lambda\asymp b^{\lambda}\lambda^{-1}$. Moreover, if
$\ell\mid b$ then
\begin{enumerate}
\item[\textup{(i)}] $\frac{1}{\pi_\lambda}\#\bigl\{p\in\mathcal P_\lambda:\ \ell\mid R_\lambda(p)\bigr\}
\ \longrightarrow\ \frac{1}{b-1}\Bigl\lfloor\frac{b-1}{\ell}\Bigr\rfloor\qquad(\lambda\to\infty),$

\item[\textup{(ii)}] $\displaystyle\sum_{m\ge2}\frac{1}{\pi_\lambda}
\#\bigl\{p\in\mathcal P_\lambda:\ \ell^{m}\mid R_\lambda(p)\bigr\}\ \ll\ 1$.
\end{enumerate}
\end{lemma}

\begin{proof}
For fixed $i$ the interval $[ib^{\lambda-1},(i+1)b^{\lambda-1})$ has length comparable to its
endpoints, so the prime number theorem in the form $\pi(z)\sim z/\log z$ already gives
\[
\pi_{\lambda,i}=\pi\bigl((i+1)b^{\lambda-1}\bigr)-\pi\bigl(ib^{\lambda-1}\bigr)
\ \sim\ \frac{b^{\lambda-1}}{(\lambda-1)\log b},
\]
Here the ratio $(i+1)/i$ of the endpoints is bounded away from $1$ in terms of $b$ alone, so the
difference of the two counts is comparable to each of them and the $o(1)$ relative errors are
absorbed. Summing over
$i$ gives $\pi_\lambda\sim(b-1)b^{\lambda-1}/((\lambda-1)\log b)$, whence
$\pi_{\lambda,i}\sim\pi_\lambda/(b-1)$: the leading digit of a prime in $\mathcal P_\lambda$ is
asymptotically uniformly distributed on $\{1,\dots,b-1\}$.

For (i), Lemma~\ref{lem:cong}(ii) shows that for $\ell\mid b$ one has $\ell\mid R_\lambda(p)$ if and
only if $\ell$ divides the leading digit $\eps_{\lambda-1}(p)$; exactly
$\lfloor(b-1)/\ell\rfloor$ of the possible leading digits $i\in\{1,\dots,b-1\}$ qualify, so summing
$\pi_{\lambda,i}\sim\pi_\lambda/(b-1)$ over these $i$ gives (i).

For (ii), write $v:=v_\ell(b)\ge1$ and, for $m\ge1$, put $j:=\lceil m/v\rceil$, so that $\ell^{m}\mid
b^{j}$. Suppose first that $j\le\lambda/2$. By \eqref{eq:defR},
$R_\lambda(n)\equiv\sum_{r<j}\eps_{\lambda-1-r}(n)b^{r}\pmod{\ell^{m}}$, so the condition
$\ell^{m}\mid R_\lambda(n)$ depends only on the $j$ leading digits of $n$; and as
$(\eps_{\lambda-1},\dots,\eps_{\lambda-j})$ runs through $\{0,\dots,b-1\}^{j}$ the quantity
$\sum_{r<j}\eps_{\lambda-1-r}b^{r}$ runs bijectively through $\{0,\dots,b^{j}-1\}$, so the number of
admissible leading blocks is at most $b^{j}\ell^{-m}+1$. Each fixed leading block confines $p$ to an
interval of length $b^{\lambda-j}\ge b^{\lambda/2}$, in which the Brun--Titchmarsh inequality
\cite[Thm.~2]{MV} gives $\ll b^{\lambda-j}/\bigl((\lambda-j)\log b\bigr)\ll\pi_\lambda b^{-j}$ primes,
since $(\lambda-j)^{-1}\le2\lambda^{-1}$ and $\pi_\lambda\asymp b^{\lambda}\lambda^{-1}$. Hence
$\#\{p\in\mathcal P_\lambda:\ell^{m}\mid R_\lambda(p)\}\ll\pi_\lambda(\ell^{-m}+b^{-j})$, and summing over
those $m\ge2$ with $\lceil m/v\rceil\le\lambda/2$ contributes
$\ll\sum_{m\ge2}\ell^{-m}+v\sum_{j\ge1}b^{-j}\ll1$.

Suppose now that $j>\lambda/2$, so that $m>v(\lambda/2-1)\ge\lambda/2-1$. Since $R_\lambda$ is injective
on $\mathcal P_\lambda$, the values $R_\lambda(p)$ with $p\in\mathcal P_\lambda$ are distinct positive
integers below $b^{\lambda}$, so that
\begin{equation}\label{eq:trivbound}
\#\bigl\{p\in\mathcal P_\lambda:\ d\mid R_\lambda(p)\bigr\}
\ \le\ \#\bigl\{1\le n<b^{\lambda}:\ d\mid n\bigr\}\ \le\ b^{\lambda}/d
\qquad(d\ge1),
\end{equation}
so, using $\pi_\lambda\asymp b^{\lambda}\lambda^{-1}$ and $\ell\ge2$, the remaining $m$ contribute
$\ll\lambda\sum_{m>\lambda/2-1}\ell^{-m}\ll\lambda\,2^{-\lambda/2+1}\ll1$.
\end{proof}

\begin{example}\label{ex:base10}
For $b=10$ we have $b(b^{2}-1)=2\cdot3^{2}\cdot5\cdot11$, so $\mathcal E_{10}=\{2,3,5,11\}$. By
Lemma~\ref{lem:cong}(i), $3\nmid R_\lambda(p)$ and $11\nmid R_\lambda(p)$ for every prime $p>11$, while
by Lemma~\ref{lem:dens}(i) the densities of the events $2\mid R_\lambda(p)$ and $5\mid R_\lambda(p)$
tend to $4/9$ and $1/9$ rather than to $1/2$ and $1/5$.
\end{example}

\subsection{Level of distribution and proof of the theorems}

Write $\pi_\lambda(z)$ for the number of $p$ with $b^{\lambda-1}\le p<z$, and
$\overline\pi_\lambda(z,a,d)$ for the number of those $p$ with $R_\lambda(p)\equiv a\pmod d$, residues
being taken in $\{1,\dots,d\}$, so that $a=d$ is the class $0\bmod d$; thus
$\pi_\lambda=\pi_\lambda(b^{\lambda})$.

\begin{theorem}[Dartyge, Rivat and Swaenepoel {\cite[Thm.~1.3]{DRS}}]\label{thm:drs}
For every integer $b\ge2$ there is an explicit constant $\xi_0=\xi_0(b)>0$ such that, given
$\xi\in(0,\xi_0)$, there exist $c=c(b,\xi)>0$ and $\lambda_0(b,\xi)\ge2$ with the following property. For
every integer $\lambda\ge\lambda_0(b,\xi)$,
\[
\sum_{\substack{d\le b^{\xi\lambda}\\ \gcd(d,\,b(b^{2}-1))=1}}\ \
\sup_{b^{\lambda-1}\le z\le b^{\lambda}}\ \ \sup_{1\le a\le d}
\left|\ \overline\pi_\lambda(z,a,d)-\frac{\pi_\lambda(z)}{d}\ \right|
\ \ll\ b^{\lambda-c\sqrt{\lambda}} .
\]
\end{theorem}

The supremum over $z$ is what permits the restriction to $\mathcal P_{\lambda,i}$ in the next lemma:
for $\mathcal A_\lambda$ itself the single value $z=b^{\lambda}$ would serve, so the localised form of
Theorem~\ref{thm:main} is bought precisely by that supremum.

\begin{lemma}\label{lem:hyp2}
Let $\xi\in(0,\xi_0(b))$ and let $i\in\{1,\dots,b-1\}$. Then for every $A>0$,
\[
\sum_{\substack{d\le b^{\xi\lambda}\\ d\in\mathcal D(\mathcal E_b)}}\bigl|r_d(\lambda)\bigr|
\ \ll_A\ \pi_{\lambda,i}\,\lambda^{-A},
\]
where $r_d(\lambda)$ is formed from $\mathcal A_{\lambda,i}$ as in \eqref{eq:defr}, and the sum is over all
such $d$ and not only the squarefree ones.
\end{lemma}

\begin{proof}
Since $\mathcal P_{\lambda,i}$
consists of the $p$ with $ib^{\lambda-1}\le p<(i+1)b^{\lambda-1}$,
\[
\#\bigl\{n\in\mathcal A_{\lambda,i}:\ d\mid n\bigr\}
=\overline\pi_\lambda\bigl((i+1)b^{\lambda-1},d,d\bigr)-\overline\pi_\lambda\bigl(ib^{\lambda-1},d,d\bigr),
\]
\[ \pi_{\lambda,i}=\pi_\lambda\bigl((i+1)b^{\lambda-1}\bigr)-\pi_\lambda\bigl(ib^{\lambda-1}\bigr),
\]
both endpoints lying in $[b^{\lambda-1},b^{\lambda}]$. Subtracting, the main terms combine to give exactly
$\pi_{\lambda,i}/d$, so $|r_d(\lambda)|$ is at most twice the supremum appearing in
Theorem~\ref{thm:drs}, and that theorem gives $\sum_{d}|r_d(\lambda)|\ll b^{\lambda-c\sqrt\lambda}$ with
$c=c(b,\xi)>0$. By Lemma~\ref{lem:dens}, $\pi_{\lambda,i}\asymp b^{\lambda}\lambda^{-1}$, so the
right-hand side is $\pi_{\lambda,i}\cdot O(\lambda b^{-c\sqrt\lambda})$; and
$b^{-c\sqrt\lambda}=\exp(-c\sqrt\lambda\,\log b)$ decays faster than any fixed
power of $\lambda^{-1}$.
\end{proof}

\begin{lemma}\label{lem:hyp3}
For every $i\in\{1,\dots,b-1\}$ one has
$\sum_{p\in\mathcal P_{\lambda,i}}\bigl(\Omega(R_\lambda(p))-\omega(R_\lambda(p))\bigr)\ll\pi_{\lambda,i}$.
\end{lemma}

\begin{proof}
Since $\mathcal P_{\lambda,i}\subseteq\mathcal P_\lambda$ and $\pi_{\lambda,i}\asymp\pi_\lambda$ by
Lemma~\ref{lem:dens}, it suffices to prove the bound with $\mathcal P_\lambda$ and $\pi_\lambda$ in place
of $\mathcal P_{\lambda,i}$ and $\pi_{\lambda,i}$; and by the identity
$\Omega(n)-\omega(n)=\sum_{m\ge2}\#\{\ell:\ \ell^{m}\mid n\}$ used in the proof of
Lemma~\ref{lem:palhyp} it suffices to show that
\[
\sum_{m\ge2}\ \sum_{\ell}\ \frac{1}{\pi_\lambda}
\#\bigl\{p\in\mathcal P_\lambda:\ \ell^{m}\mid R_\lambda(p)\bigr\}\ \ll\ 1 .
\]
Fix $\xi\in(0,\xi_0(b))$ and split the range of $(\ell,m)$ into four parts. Every prime power counted
divides some $R_\lambda(p)<b^{\lambda}$, so $2\le m\le M:=\lambda\log b/\log2$ throughout.

If $\ell\mid b^{2}-1$ then $\ell\nmid R_\lambda(p)$ for $p>b+1$ by Lemma~\ref{lem:cong}(i), so these
$\ell$ contribute $O(1)$ in total; and if $\ell\mid b$ then Lemma~\ref{lem:dens}(ii) bounds the
corresponding subsum by $O(1)$ for each of the $O(1)$ such $\ell$.

Let now $\ell\nmid b(b^{2}-1)$. If $\ell^{m}\le b^{\xi\lambda}$ then $\ell^{m}\in\mathcal D(\mathcal E_b)$,
so Lemma~\ref{lem:hyp2} applies. Distinct pairs $(\ell,m)$ with $m\ge2$ give distinct moduli $\ell^{m}$, so
summing the conclusion of that lemma over the $b-1$ values of $i$ and using
$\#\{p\in\mathcal P_\lambda:\ell^{m}\mid R_\lambda(p)\}=\pi_\lambda\ell^{-m}+\sum_{i}r^{(i)}_{\ell^{m}}(\lambda)$,
where $r^{(i)}$ is formed from $\mathcal A_{\lambda,i}$, we obtain
\[
\sum_{m\ge2}\ \sum_{\substack{\ell\nmid b(b^{2}-1)\\ \ell^{m}\le b^{\xi\lambda}}}
\frac{1}{\pi_\lambda}\#\bigl\{p:\ \ell^{m}\mid R_\lambda(p)\bigr\}
\ \le\ \sum_{m\ge2}\sum_{\ell}\frac{1}{\ell^{m}}+O\bigl(\lambda^{-1}\bigr)\ =\ O(1).
\]
Finally, still with $\ell\nmid b(b^{2}-1)$, if $\ell^{m}>b^{\xi\lambda}$ then
$\ell>u_m:=b^{\xi\lambda/m}$, and \eqref{eq:trivbound} with
$\pi_\lambda\asymp b^{\lambda}\lambda^{-1}$ and
$\sum_{n>u}n^{-m}\le u^{-m}+u^{1-m}/(m-1)\le2u^{1-m}$, valid for $u\ge1$ and $m\ge2$, gives
\[
\begin{aligned}
\sum_{2\le m\le M}\ \sum_{\ell>u_m}\frac{1}{\pi_\lambda}
\#\bigl\{p:\ \ell^{m}\mid R_\lambda(p)\bigr\}
&\ \ll\ \lambda\sum_{2\le m\le M}\ \sum_{\ell>u_m}\frac{1}{\ell^{m}} \\
&\ \ll\ \lambda\sum_{2\le m\le M}u_m^{1-m}\ \ll\
\lambda^{2}b^{-\xi\lambda/2}\ \ll\ 1,
\end{aligned}
\]
since $u_m^{1-m}=b^{-\xi\lambda(m-1)/m}\le b^{-\xi\lambda/2}$ for $m\ge2$ and the sum over $m$ has
$O(\lambda)$ terms.
\end{proof}

\begin{proof}[Proof of Theorem~\ref{thm:main}, and of Theorem~\ref{thm:moments} for reversed primes]
Fix $i\in\{1,\dots,b-1\}$ and apply Proposition~\ref{prop:criterion} to
$\mathcal B_\lambda:=\mathcal A_{\lambda,i}$ with $\mathcal E:=\mathcal E_b$, $\xi\in(0,\xi_0(b))$
fixed, and $S$ the set of all primes.\footnote{By Lemma~\ref{lem:dens} we have $\pi_{\lambda,i}>0$ for
all large $\lambda$; we set $\mathcal B_\lambda:=\{1\}$ for the finitely many $\lambda$ at which
$\mathcal A_{\lambda,i}$ is empty. This does not affect the limits, since all statements are asymptotic
in $\lambda$.} Hypothesis (i) holds by
\eqref{eq:defR}, which gives $R_\lambda(p)<b^{\lambda}$ for every $p\in\mathcal P_\lambda$; hypothesis
(ii) is Lemma~\ref{lem:hyp2}, and hypothesis (iii) is
Lemma~\ref{lem:hyp3}. This gives Theorem~\ref{thm:moments} for $\mathcal A_{\lambda,i}$ and the localised
form of Theorem~\ref{thm:main}.

Since $\mathcal P_\lambda$ is the disjoint union of the $\mathcal P_{\lambda,i}$, every average over
$\mathcal A_\lambda$ is the convex combination, with weights $\pi_{\lambda,i}/\pi_\lambda$, of the
corresponding averages over the $\mathcal A_{\lambda,i}$. Applying this to $(\omega(n)-L)^{k}$ gives
Theorem~\ref{thm:moments} for $\mathcal A_\lambda$, the $b-1$ estimates having the same main term and the
same relative error when $k$ is even, and the same upper bound when $k$ is odd; and applying it to the
distribution functions of the localised statement,
$F_{\lambda,i}(t):=\pi_{\lambda,i}^{-1}\,\#\bigl\{p\in\mathcal P_{\lambda,i}:\ \omega(R_\lambda(p))\le
L+t\sqrt L\bigr\}$,
gives $\sup_t|F_\lambda-\Phi|\le\max_i\sup_t|F_{\lambda,i}-\Phi|=o(1)$ for the corresponding function
$F_\lambda$ formed with $\mathcal P_\lambda$ and $\pi_\lambda$, which is the unlocalised form of
Theorem~\ref{thm:main}. The argument for $\Omega$ is identical.
\end{proof}

\begin{proof}[Proof of Corollaries~\ref{cor:normalorder} and \ref{cor:cheb}]
For Corollary~\ref{cor:normalorder}, apply Theorems~\ref{thm:pal} and \ref{thm:main} with
$t=\pm\theta\sqrt L$, where $L\to\infty$, and use $\Phi(-u)+1-\Phi(u)\to0$ as $u\to\infty$. For the
last assertion, every element of $\mathcal T_\lambda\cup\mathcal A_\lambda$ lies in
$[b^{\lambda-1},b^{\lambda})$ for $\lambda\ge3$: for $\mathcal T_\lambda$ by definition, and for
$\mathcal A_\lambda$ by Lemma~\ref{lem:cong}(i), since $p\ge b^{\lambda-1}\ge b^{2}>b$ for every
$p\in\mathcal P_\lambda$. On that range $\log\log n=L+O(\lambda^{-1})$, as was already noted after
\eqref{eq:defCk}. For
Corollary~\ref{cor:cheb}, let $S:=S_{\mathcal C}$ be the set of primes unramified in $K$ whose Frobenius
class is $\mathcal C$; by the Chebotarev density theorem with the classical
error term \cite[Thm.~1.3]{LO} and partial summation,
$\sum_{\ell\le y,\ \ell\in S}\ell^{-1}=\delta\log\log y+O(1)$, so $S$ is regular of density $\delta$, and
the proofs above apply verbatim with this $S$.
\end{proof}

\section{Remarks}\label{sec:remarks}

\begin{remark}\label{rem:inputs}
The proofs require a level of distribution with moduli up to a fixed power of $b^{\lambda}$, together with
the treatment of the primes dividing $b(b^{2}-1)$ in \eqref{eq:defE} and Lemma~\ref{lem:dens}. Any
improvement of $\xi_0(b)$ in Theorem~\ref{thm:drs}, or of $\beta(b)$ in Theorem~\ref{thm:col}, leaves the
results unchanged. A Siegel--Walfisz range for the moduli would not suffice: for reversed primes the range
$d\le\exp(c\sqrt\lambda)$ of \cite{BS1,BS2}, and for palindromes the range
$q\le\exp(c\lambda/\log\lambda)$ of \cite[Thm.~1]{Col}, would force the truncation point $y$ of
\eqref{eq:defy} below every fixed power of $b^{\lambda}$, and the bound
$\#\{\ell>y:\ell\mid n\}\le\lambda\log b/\log y$ of Lemma~\ref{lem:setup}(ii) would then exceed
$\sqrt{L}$.
\end{remark}

\begin{remark}\label{rem:extremal}
Theorem~\ref{thm:pal} determines the typical value of $\omega$ on $\mathcal T_\lambda$ and says nothing
about its largest value, which is the question of \cite{BSh} and remains open: the trivial upper bound is
$\max_{n\in\mathcal T_\lambda}\omega(n)\ll\lambda/\log\lambda$, while the lower bound of \cite{BSh} is
$(\log\log n)^{1+o(1)}=\lambda^{o(1)}$. The two are far apart, and Theorem~\ref{thm:pal} does not narrow
the gap. It does recover the lower bound in the form in which it is stated, and for a positive proportion
of $n\in\mathcal T_\lambda$ rather than for a single one: for each fixed $t$ the proportion of
$n\in\mathcal T_\lambda$ with $\omega(n)\ge\log\log n+t\sqrt{\log\log n}$ tends to $1-\Phi(t)>0$, and
$\log\log n+t\sqrt{\log\log n}=(\log\log n)^{1+o(1)}$.
\end{remark}

\begin{remark}\label{rem:paldist}
Tuxanidy and Panario \cite[Thm.~1.5]{TP} prove a level of distribution $1/5-\eta$ for palindromes, with an
exponent independent of $b$, whereas Col's $\beta(b)\sim1/(6\pi b)$ degrades as $b$ grows. Their result is
stated for the palindromes coprime to $b^{3}-b$, and every palindrome with an even number of digits is
divisible by $b+1$; the set to which it applies therefore contains no element of $\mathcal T_\lambda$ when
$\lambda$ is even. Substituting it for Theorem~\ref{thm:col} in Lemma~\ref{lem:palhyp} would give
Theorem~\ref{thm:pal} for the palindromes in $\mathcal T_\lambda$ coprime to $b^{3}-b$, with $\xi$ close
to $1/5$, but only for $\lambda$ odd. For the squarefree palindromes an exponent of distribution $1/16$
is due to Tuxanidy \cite{Tux}.
\end{remark}

\begin{remark}\label{rem:rate}
Theorem~\ref{thm:moments} is quantitative, but we do not deduce from it a rate of convergence in
Theorem~\ref{thm:pal} or Theorem~\ref{thm:main}. The classical method of moments \cite[Thm.~30.2]{Bil} is
qualitative, and the known routes to a rate for $\omega$ proceed instead through characteristic functions,
as in R\'enyi and Tur\'an \cite{RT}, or through Stein's method, as in Harper \cite{Har}. For a digitally
defined sequence a rate has been obtained by the former route: Mauduit and S\'ark\"ozy
\cite[Thm.~4]{MS2} give $\ll(\log\log\log N)/\sqrt{\log\log N}$ for the integers up to $N$ with a fixed
digit sum. A cumulant
criterion such as \cite[Thm.~2.4]{DJS} does convert into a Kolmogorov bound, but requires control of
cumulants of every order.
\end{remark}

\begin{remark}[Formal verification]\label{rem:lean}
The deductive content of this note has been verified with the Lean~4 proof
assistant (Lean~4.33.0 with Mathlib). In the companion repository
\begin{center}
  \url{https://github.com/vibefrtz/vibemath}\\[2pt]
  (folder \texttt{erdos-kac-reversed-primes}, archived at
  \href{https://doi.org/10.5281/zenodo.22078540}{\texttt{doi:10.5281/zenodo.22078540}})
\end{center}
the results quoted from the literature --- Proposition~\ref{prop:gs},
Theorems~\ref{thm:col}, \ref{thm:bsh} and \ref{thm:drs}, the
Brun--Titchmarsh inequality, the prime number theorem, Mertens' second
theorem and the method of moments --- are declared as axioms, in the form in
which they are used here, and everything this note itself proves is
machine-checked from those axioms with no unproved goals:
Proposition~\ref{prop:criterion} together with Lemma~\ref{lem:setup} and the
moment computation; Lemmas~\ref{lem:palhyp}, \ref{lem:cong}, \ref{lem:dens},
\ref{lem:hyp2} and \ref{lem:hyp3}; and the deductions of
Theorems~\ref{thm:pal} and \ref{thm:main} (pointwise and uniformly in $t$,
localised and unlocalised), of Theorem~\ref{thm:moments} at each fixed
moment order, and of Corollaries~\ref{cor:normalorder} and \ref{cor:cheb},
the latter for an arbitrary regular set $S$. An axiom audit included in the
repository lists, for every theorem, the exact axioms its proof consumes;
these coincide with the citations invoked by the corresponding proof in this
note. Two points lie outside the verification and are documented in the
repository: the transcriptions of the quoted results were not checked
mechanically against their sources beyond what is recorded there, and the
uniformity in $k$ asserted in Theorem~\ref{thm:moments} is verified at each
fixed $k$ only --- the case that Theorems~\ref{thm:pal} and \ref{thm:main}
use.
\end{remark}


\section*{Declaration on the use of AI tools}

This note was prepared with extensive assistance from large language models
(Claude Fable 5 and Claude Opus 5, Anthropic), used at every stage of its preparation: for the
literature search, for jointly working out the main arguments, and for drafting and typesetting the
manuscript itself. In particular,
Theorem~\ref{thm:col}, on which Theorem~\ref{thm:pal} depends, was located by
the models, and the modification of \cite[Prop.~3]{GS} presented in
Proposition~\ref{prop:criterion} was worked out jointly with the models. The
author has verified all statements, proofs and references, made the final
decisions on content and presentation, and is responsible for any errors that
remain.

\begin{thebibliography}{99}

\bibitem{BSh} W.~D.~Banks and I.~E.~Shparlinski, \emph{Prime divisors of palindromes},
Period.\ Math.\ Hungar.\ \textbf{51} (2005), 1--10.

\bibitem{BS1} G.~Bhowmik and Y.~Suzuki, \emph{On Telhcirid's theorem on arithmetic progressions},
preprint, 2024, \href{https://arxiv.org/abs/2406.13334}{arXiv:2406.13334}.

\bibitem{BS2} G.~Bhowmik and Y.~Suzuki, \emph{The Zsiflaw--Legeis theorem for arbitrary bases},
preprint, 2025, \href{https://arxiv.org/abs/2507.08714}{arXiv:2507.08714}.

\bibitem{Bil} P.~Billingsley, \emph{Probability and measure}, 3rd ed., Wiley, New York, 1995.

\bibitem{CJ} S.~Chourasiya and D.~R.~Johnston, \emph{Power-free palindromes and reversed primes},
Monatsh.\ Math.\ \textbf{208} (2025), no.~3, 369--385.

\bibitem{Col} S.~Col, \emph{Palindromes dans les progressions arithm\'etiques},
Acta Arith.\ \textbf{137} (2009), 1--41.

\bibitem{DMRSS} C.~Dartyge, B.~Martin, J.~Rivat, I.~E.~Shparlinski and C.~Swaenepoel,
\emph{Reversible primes},
J.\ Lond.\ Math.\ Soc.\ (2) \textbf{109} (2024).

\bibitem{DRS} C.~Dartyge, J.~Rivat and C.~Swaenepoel, \emph{Prime numbers with an almost prime reverse},
preprint, 2025, \href{https://arxiv.org/abs/2506.21642}{arXiv:2506.21642}.

\bibitem{DJS} H.~D\"oring, S.~Jansen and K.~Schubert, \emph{The method of cumulants for the normal
approximation}, Probab.\ Surv.\ \textbf{19} (2022), 185--270.

\bibitem{EK} P.~Erd\H os and M.~Kac, \emph{The Gaussian law of errors in the theory of additive number
theoretic functions}, Amer.\ J.\ Math.\ \textbf{62} (1940), 738--742.

\bibitem{EMS} P.~Erd\H os, C.~Mauduit and A.~S\'ark\"ozy, \emph{On arithmetic properties of integers
with missing digits. II. Prime factors}, Discrete Math.\ \textbf{200} (1999), 149--164.

\bibitem{GS} A.~Granville and K.~Soundararajan, \emph{Sieving and the Erd\H os--Kac theorem}, in:
Equidistribution in number theory, an introduction, A.~Granville and Z.~Rudnick (eds.),
NATO Sci.\ Ser.\ II Math.\ Phys.\ Chem.\ \textbf{237}, Springer, Dordrecht, 2007, pp.\ 15--27.

\bibitem{HJ} M.~Harm and D.~R.~Johnston, \emph{The reverse Goldbach problem and a refined
Zsiflaw--Legeis theorem}, preprint, 2026, \href{https://arxiv.org/abs/2605.21876}{arXiv:2605.21876}.

\bibitem{Har} A.~J.~Harper, \emph{Two new proofs of the Erd\H os--Kac theorem, with bound on the rate of
convergence, by Stein's method for distributional approximations}, Math.\ Proc.\ Cambridge Philos.\ Soc.\
\textbf{147} (2009), 95--114.

\bibitem{JK} D.~R.~Johnston and B.~Kerr, \emph{The infinitude of square-free palindromes},
preprint, 2026, \href{https://arxiv.org/abs/2601.07097}{arXiv:2601.07097}.

\bibitem{KMS} S.~Konyagin, C.~Mauduit and A.~S\'ark\"ozy, \emph{On the number of prime factors of
integers characterized by digit properties}, Period.\ Math.\ Hungar.\ \textbf{40} (2000), 37--52.

\bibitem{LO} J.~C.~Lagarias and A.~M.~Odlyzko, \emph{Effective versions of the Chebotarev density
theorem}, in: Algebraic number fields: $L$-functions and Galois properties, A.~Fr\"ohlich (ed.),
Academic Press, London, 1977,
pp.\ 409--464.

\bibitem{MS1} C.~Mauduit and A.~S\'ark\"ozy, \emph{On the arithmetic structure of sets characterized
by sum of digits properties}, J.\ Number Theory \textbf{61} (1996), 25--38.

\bibitem{MS2} C.~Mauduit and A.~S\'ark\"ozy, \emph{On the arithmetic structure of the integers whose
sum of digits is fixed}, Acta Arith.\ \textbf{81} (1997), 145--173.

\bibitem{MV} H.~L.~Montgomery and R.~C.~Vaughan, \emph{The large sieve},
Mathematika \textbf{20} (1973), 119--134.

\bibitem{MMP} M.~R.~Murty, V.~K.~Murty and S.~Pujahari, \emph{An all-purpose Erd\H os--Kac theorem},
Math.\ Z.\ \textbf{305} (2023).

\bibitem{RT} A.~R\'enyi and P.~Tur\'an, \emph{On a theorem of Erd\H os--Kac}, Acta Arith.\ \textbf{4}
(1958), 71--84.

\bibitem{Tux} A.~Tuxanidy, \emph{Distribution of square-free palindromes}, preprint, 2026,
\href{https://arxiv.org/abs/2603.02549}{arXiv:2603.02549}.

\bibitem{TP} A.~Tuxanidy and D.~Panario, \emph{Infinitude of palindromic almost-prime numbers},
Int.\ Math.\ Res.\ Not.\ IMRN \textbf{2024}.

\end{thebibliography}

\end{document}
